% =====================================================================
%  DRAFT MANUSCRIPT v0.6. INCOMPLETE. NOT FOR CIRCULATION.
%  Autonomous Driving Agent / CARLA 0.9.16
%  Compile on Overleaf with pdfLaTeX. Class: IEEEtran (conference).
%
%  STATUS OF THIS FILE
%    Sections I to V : drafted
%    Section  VI     : RENDERING FAILURES, drafted but needs measured numbers
%    Section  VII    : RESULTS, deliberately empty, corpus does not exist
%    Section  VIII   : DISCUSSION, stub only
%    Section  IX     : CONCLUSION, unwritten
%    References      : partial, renumber after additions
%  Search for "TODO" to find every open item.
%  To strip all draft markers: delete the \TODO/\TODOflow/\UNWRITTEN calls
%  and their definitions below, plus the red banner in \title.
%
%  NOTE ON HYPHENS: automatic hyphenation at line breaks is switched off
%  below, and compound hyphens have been removed from the prose. The only
%  hyphens left in the file are inside quoted paper titles in the
%  bibliography, which must stay as published.
% =====================================================================

\documentclass[conference]{IEEEtran}
\IEEEoverridecommandlockouts

\usepackage{cite}
\usepackage{amsmath,amssymb}
\usepackage{graphicx}
\usepackage{booktabs}
\usepackage{array}
\newcolumntype{L}[1]{>{\raggedright\arraybackslash}p{#1}}
\usepackage{xcolor}
\usepackage{textcomp}
\usepackage{framed}
\usepackage[hidelinks]{hyperref}

% No hyphenation at line breaks anywhere in the document.
\hyphenpenalty=10000
\exhyphenpenalty=10000
\tolerance=9999
\emergencystretch=3em
\sloppy
\widowpenalty=0
\clubpenalty=0

% Markers for the unfinished parts.
\colorlet{shadecolor}{white}
\newcommand{\TODO}[1]{%
  \begingroup
  \setlength{\FrameSep}{4pt}%
  \def\FrameCommand{{\color{gray!55}\vrule width 1.2pt}\hspace{5pt}}%
  \MakeFramed{\advance\hsize-\width \FrameRestore}%
  \noindent\scriptsize\itshape\color{gray!80}#1%
  \endMakeFramed
  \endgroup}

\newcommand{\UNWRITTEN}[1]{\par\noindent{\itshape\color{gray!85}#1}\par\vspace{2pt}}

% Same kind of note, but breakable across a column (a framed box is not).
\newcommand{\TODOflow}[1]{%
  \par\noindent{\scriptsize\itshape\color{gray!80}#1\par}}

\begin{document}

\title{Active Imitation Learning for a Neural Network\\ That Drives a Car in CARLA\\[3pt]
{\normalsize\color{red!65!black}\bfseries DRAFT v0.6. INCOMPLETE MANUSCRIPT. NOT FOR CIRCULATION.}}

\author{
\IEEEauthorblockN{
Medhansh Khattar,\quad
Prashant Kumar Jaglan,\quad
Yash Yadav\\[2pt]
Samyak Jain,\quad
Sanchit Kumar Sinha,\quad
Aryan Garg}
\IEEEauthorblockA{Vellore Institute of Technology, Bhopal, India}
}

\maketitle

% =====================================================================
\begin{abstract}
We describe an imitation learning pipeline for a neural network that
drives a car in CARLA 0.9.16. The policy is a small multilayer
perceptron of 667 parameters that reads eight scalars describing the car
relative to its route and outputs steering, accelerator and brake. Three
contracts hold the system together: an observation contract of eight
fields, an action protocol of seven fields, and a deployability contract
that an exported model must satisfy before the runtime will let it
drive. A safety supervisor runs outside the policy on the simulator host
and keeps full braking authority.

The title names our aim, not our current state: active imitation, where
the learner drives and the expert corrects the states it visits, is the
next milestone, and what runs today is passive cloning. The method is
complete, but the evaluation is not. We cannot yet say how well any of
it drives, because we cannot keep the simulator rendering long enough to
record one complete route. Camera compositing saturates the integrated
GPU at speed, environment streaming stalls the client, and the crashes
that follow end episodes partway through. Our acceptance rule discards
any interrupted route, so these crashes show up as an empty training
set: zero accepted routes against a target of twelve.
\textit{Sections VII to IX are left unwritten.}
\end{abstract}

\begin{IEEEkeywords}
autonomous driving, active imitation learning, behaviour cloning, CARLA,
rendering bottleneck, safety supervision, reproducibility.
\end{IEEEkeywords}

% =====================================================================
\section{Introduction}

CARLA~\cite{carla} is the most widely used open simulator for driving
research. It runs an Unreal Engine world behind an RPC interface,
supplies OpenDRIVE road topology, traffic and weather, and replays a
seeded route deterministically after a code change. Crashes cost
nothing, which is the usual argument for working in simulation. That
argument assumes the simulator keeps running.

Most work since Bojarski \textit{et al.}~\cite{bojarski} maps camera
pixels straight to control. We started somewhere narrower. Our policy
reads eight scalars measured relative to the route instead of images,
mainly because of how few demonstrations six undergraduates can curate:
on a few thousand frames, a larger network learns the shape of the
training town rather than how to stay in a lane~\cite{limits}. What
follows is the architecture, the contracts holding it together, and the
rendering failures we spent most of the project fighting. Where a
driving result would normally go, we say so, and Section~VI explains
why.

% =====================================================================
\section{Related Work}

Behaviour cloning for driving was shown at scale by Bojarski \textit{et
al.}~\cite{bojarski}. Codevilla \textit{et al.} added a command input so
one network could resolve what to do at an intersection~\cite{cil}, and
later catalogued what cloning gets wrong on its own~\cite{limits}:
distributional shift, and a bias toward whatever action the expert takes
most often. Ross \textit{et al.} give the usual fix,
DAgger~\cite{dagger}, in which the expert labels the states the learner
visits. That is the active imitation of our title, and the step we have
not reached. Within CARLA, Learning by Cheating~\cite{lbc} beats direct
sensorimotor cloning by distilling from a privileged agent, and
TransFuser~\cite{transfuser} fuses several sensors. Our observation
contract sits closer to the privileged stage of~\cite{lbc}.

\TODO{TODO (Related Work). Still to read and cite: Toromanoff
\textit{et al.} on implicit affordances (CVPR 2020); the Kiran
\textit{et al.} DRL survey (2021), for positioning; SAE J3016, so the
autonomy claim is precise; and anything on simulator throughput, which
we have not looked at and which bears on Section~VI.}

% =====================================================================
\section{System Architecture}

Two machines share the work over a direct Ethernet link. Machine~1 runs
the CARLA server, which owns physics, the OpenDRIVE map, the traffic
manager, sensors and the vehicle. Machine~2 runs inference, the route
planner, an observer window and network visualisation. Each iteration
goes: machine~1 publishes world state, machine~2 runs the planner and
the policy, machine~2 sends back a driving command, and machine~1's
safety layer checks that command before applying it.

We split the map in two. The physical map stays with the simulator,
which owns the world. The navigation map needs only road topology to
pick a destination and lay out waypoints, so it can live on machine~2. A
watchdog on machine~1 is meant to brake if the link drops. It is
written, but we have not tested it by pulling the cable.

% =====================================================================
\section{Methodology}

\subsection{Observation contract}

\texttt{policy\_contract.py} defines what the network sees: the eight
scalars in Table~\ref{tab:obs}, each clamped by a normaliser the same
module owns.

\begin{table}[htbp]
\caption{The observation contract, eight fields}
\label{tab:obs}
\centering\scriptsize
\setlength{\tabcolsep}{3pt}
\begin{tabular}{@{}L{2.25cm}L{1.4cm}L{2.0cm}L{1.45cm}@{}}
\toprule
Field & Meaning & Field & Meaning \\
\midrule
\texttt{lane\_error\_m}     & offset from lane centre & \texttt{front\_obst\_m} & nearest car ahead \\
\texttt{heading\_err\_deg}  & car versus road angle & \texttt{red\_light}     & binary: light governs \\
\texttt{speed\_kmh}         & current road speed      & \texttt{curvature}      & bend $\approx$8\,m ahead \\
\texttt{target\_spd\_kmh}   & speed controller wants  & \texttt{stopped}        & binary: $<$1\,km/h \\
\bottomrule
\end{tabular}
\end{table}

Working in coordinates relative to the route rather than to the world is
what lets a policy trained on one town be evaluated on another. One gap
is still open. The waypoint used as the route reference when computing
heading error and curvature during training is not handed to the policy
the same way at inference, so the two can disagree on those two fields.
Any number we report before closing this would be suspect.

\subsection{Action protocol and safety supervision}

The policy produces three values and the vehicle interface takes seven.
Steering comes out in $[-1,+1]$ and maps to $\pm70^{\circ}$. Accelerator
and brake are clamped to $[0,1]$, and the runtime will not apply both at
once. The output layer is linear, so all clamping happens in the runtime
rather than inside the network. Clutch, handbrake, wipers and gear come
from a deterministic transmission model with six gears. CARLA does not
simulate engine speed, so we estimate it from road speed and gear ratio
and label it \emph{estimated telemetry} everywhere it appears.

The safety supervisor runs on the simulator host, outside the policy. It
forces a full brake on collision, sensor timeout, a command older than
0.5\,s, lane error above 2\,m, an obstacle inside 7\,m, speed above
50\,km/h, a red light, or telemetry that fails validation. Nothing else
is allowed to override the network.

\subsection{Recording demonstrations}

We record from two experts: CARLA's \texttt{BasicAgent}, which follows
fixed rules, and a person driving the cockpit. A route runs from a spawn
point to a destination at least 150\,m away, and each tick appends one
JSONL row with the eight observations, the expert's three controls,
telemetry and scenario metadata. A route is kept only if it finished,
hit nothing, and holds at least 100 samples. Crashed and interrupted
runs are thrown out rather than trimmed. That rule matters for
Section~VI, because it is what turns rendering failures into missing
data instead of bad data.

Recording happens on Town05. Town04 is held back as an unseen map, so a
good final score cannot be explained by the network having memorised the
road. Twelve Town04 teacher episodes recorded before we adopted this
rule are quarantined, not deleted.

\subsection{Training}

Training runs offline in a Colab notebook and writes a JSON artifact the
local driver loads. The network is a linear layer from 8 to 24, $\tanh$,
a linear layer from 24 to 16, $\tanh$, then a linear head from 16 to 3,
giving 667 parameters. We use AdamW at $\mathrm{lr}=10^{-3}$ with weight
decay $10^{-5}$, mean squared error over the three control channels,
batches of 128, at most 80 epochs, and early stopping after 12 epochs
without improvement. Seed~7 is set for Python, NumPy and PyTorch, and
everything runs on CPU.

Splits are taken over whole routes, 70/15/15, never over single frames.
Neighbouring frames are nearly identical, so splitting at random puts
copies of the same moment in both train and test and the reported error
stops meaning anything. Each episode lands in one split, and an
assertion checks the three are disjoint. The notebook writes its
normaliser straight out of the contract file, so training and inference
cannot drift apart.

\subsection{Deployability gate}

Every exported artifact carries a contract version, a
\texttt{deployable} flag and the episode names in each split. The driver
rejects anything that is not contract v2, or that has the flag set
false. The only artifact we have at the moment is a synthetic smoke test
at contract v1 with \texttt{deployable} false, so it gets rejected and
route mode stays unavailable. There is no scripted fallback, so the car
simply does not drive itself. Training is disabled in the runtime, which
means a deployed model cannot quietly retrain.

One runtime behaviour is not learned, and we should declare it. Cloning
produces too much braking, so the policy tends to sit still on an empty
road. A bounded recovery capped at 0.35 accelerator releases the stall,
switches off at once for a red light, an obstacle or bad telemetry, and
never touches the weights. Any result has to say whether it was on.

% =====================================================================
\section{Experimental Setup (partial)}

The scenario catalogue names four categories, so an experiment can ask
for a category instead of fixing a map in code: urban (Town03, Town04,
Town05, Town10HD, 30 NPCs), highway (Town05, Town06, Town07, 20), rural
(Town01, Town02, 12) and intersection (Town03, Town04, 40). Weather
comes from eight presets and road damage from three pothole modes. All
of this is configuration. No category has completed a verified
evaluation run, for the reasons in the next section.

\TODO{TODO (Section V). Missing: a hardware table for both machines;
control loop latency over the link (mean, p99); headless tick rate; the
commit hash for the reported runs.}

% =====================================================================
\section{The Rendering Bottleneck}

Most of our difficulty has come from rendering rather than from the
learning algorithm. We cannot reliably keep the environment and the
camera image rendering long enough to finish a single route.

\subsection{Rendering the camera image at speed}

The visual cockpit composites four camera feeds on top of streamed
scenery. On an Intel Arc~140V integrated GPU that is too much. Frame
rate falls as road speed rises, because travelling faster makes Unreal
stream in new scenery while the same GPU is still compositing the camera
panel. It does not degrade gently. The image stalls, the client falls
behind the server tick, and the commands reaching the car go stale.

We have two mitigations, and both cost realism. The default neural path
runs \emph{headless}, with no camera feeds at all. On the visual build,
throttle is cut above 42\,km/h and top speed capped at 48\,km/h to keep
the GPU alive. The speed range we can demonstrate in is therefore set by
a rendering limit rather than by anything about driving, and we cannot
record behaviour at highway speed at all.

\subsection{Environment rendering and process failures}

Environment rendering fails in three ways, all ending the episode:

\begin{itemize}
\setlength{\itemsep}{1pt}\setlength{\topsep}{2pt}
\setlength{\parskip}{0pt}\setlength{\partopsep}{0pt}
\item \textbf{Thermal termination.} The GPU shares a thermal package
with the CPU. One afternoon of sustained rendering produced both a fatal
hardware error and a critical thermal hibernation. This hardware exposes
no live temperature counter, so our health guard reads the Windows Event
Log instead and refuses to launch during a cooldown of 30 or 60 minutes.
\item \textbf{Silent process death.} A rendering component wedges,
Windows kills it as a hung application, and the rest of the stack
follows. For a long time nothing was logged, because the old launcher
fired each process into a detached task whose output went nowhere. One
supervisor now owns both processes, writes logs, checks liveness by
probing the RPC port rather than trusting a PID, and restarts with a
backoff from 5\,s to 120\,s.
\item \textbf{Orphaned renderers.} \texttt{CarlaUE4.exe} is a shim
around another executable, and the virtualenv \texttt{python.exe} starts
the real interpreter again as a child. Killing the visible process left
the real renderer holding port 2000, so the next launch fought the
previous one. The supervisor now kills process trees and sweeps orphans
before starting.
\end{itemize}

\subsection{How this empties the training set}

These failures are why Section~VII is empty. A rendering stall
interrupts a route midway, and by the acceptance rule above an
interrupted route is discarded rather than kept, which we chose on
purpose: a route ending in a crash teaches the policy that crashing is
ordinary. The two rules do not sit well together. The more often the
environment fails to render, the fewer routes survive, and we currently
have zero accepted routes out of twelve. The network is not failing to
learn. We have not given it anything to learn from.

\TODO{TODO (Section VI). Qualitative; make quantitative before
submission. Needed: (a) FPS against road speed, headless against visual,
as a figure; (b) mean route survival time before a rendering
termination; (c) counts per failure class; (d) package temperature at
termination (unobtainable, so state it); (e) does splitting across two
machines help.}

% =====================================================================
\section{Results}

\UNWRITTEN{This section is not yet written.}

There is no corpus to write it from, for the reasons above, and a
training loss against the smoke test artifact would say nothing. To stop
ourselves picking metrics after seeing results, we are fixing the format
now: collisions per 100\,km, lane departures per km, route completion
rate, safety interventions per route, and mean absolute error for each
control channel, each for the scripted baseline and for the learned
policy on the map we held back.

\TODO{TODO (Section VII). Prerequisites, in order: (1) stabilise
rendering so a route of 150\,m survives from start to finish; (2) accept
at least 12 Town05 routes; (3) close the training and inference gap in
Section~IV; (4) export a contract v2 artifact with \texttt{deployable}
set true; (5) run the Town04 evaluation with the stall assist on and
off. No prose before (5) completes.}

% =====================================================================
\section{Discussion}

\UNWRITTEN{Stub. Only the observations that do not depend on results.}

Most of the time on this project went to rendering and process
management rather than to modelling, which is not something the papers
we started from discuss. Results in~\cite{lbc} and~\cite{transfuser}
assume a simulator that renders reliably. When it does not, the work
moves somewhere those papers have no reason to go.

\TODO{TODO (Section VIII). Once VII exists: (a) a leakage demo, error
under splitting by route against by frame; (b) a failure taxonomy from
Town04; (c) what eight scalars cannot represent (pedestrians, signs,
anything needing vision); (d) does headless collection bias the policy.}

% =====================================================================
\section{Conclusion and Future Work}

\UNWRITTEN{[Not written. To be drafted after Sections VI and VII are
complete.]}

Planned next steps, in rough order. Get the rendering stable, moving the
renderer to dedicated hardware if that is what it takes. Replace the
eight scalars with RGB, depth and semantic segmentation. Add the command
input of~\cite{cil} and the active correction loop of~\cite{dagger},
which is what would make our title accurate. Train across the whole
catalogue rather than one town. Reward logging already runs every tick,
but nothing consumes it.

% =====================================================================
% Hyphens inside these quoted titles are part of the published titles
% and must not be edited.
\begin{thebibliography}{99}
\footnotesize
\setlength{\itemsep}{0pt}\setlength{\parsep}{0pt}

\bibitem{carla} A. Dosovitskiy \textit{et al.}, ``CARLA: An Open Urban
Driving Simulator,'' \emph{CoRL}, 2017. arXiv:1711.03938.

\bibitem{bojarski} M. Bojarski \textit{et al.}, ``End to End Learning
for Self-Driving Cars,'' NVIDIA, 2016. arXiv:1604.07316.

\bibitem{cil} F. Codevilla \textit{et al.}, ``End-to-end Driving via
Conditional Imitation Learning,'' \emph{ICRA}, 2018. arXiv:1710.02410.

\bibitem{limits} F. Codevilla \textit{et al.}, ``Exploring the
Limitations of Behavior Cloning for Autonomous Driving,'' \emph{ICCV},
2019.

\bibitem{dagger} S. Ross, G. Gordon, and J. A. Bagnell, ``A Reduction of
Imitation Learning and Structured Prediction to No-Regret Online
Learning,'' \emph{AISTATS}, 2011. arXiv:1011.0686.

\bibitem{lbc} D. Chen \textit{et al.}, ``Learning by Cheating,''
\emph{CoRL}, 2019. arXiv:1912.12294.

\bibitem{transfuser} A. Prakash, K. Chitta, and A. Geiger,
``Multi-Modal Fusion Transformer for End-to-End Autonomous Driving,''
\emph{CVPR}, 2021.

\end{thebibliography}

\TODOflow{TODO (References). Incomplete: add CARLA docs, OpenDRIVE, rendering refs.}

\end{document}
